\documentclass[11pt]{article}
\usepackage[a4paper,margin=1in]{geometry}
\usepackage{fontspec}
\usepackage[boldfont=false]{xeCJK}
\setCJKmainfont{FandolSong-Regular.otf}
\usepackage{amsmath}
\usepackage{amssymb}
\usepackage{array}
\usepackage{booktabs}
\usepackage{graphicx}
\usepackage{adjustbox}
\usepackage{enumitem}
\setlist[itemize]{topsep=2pt,partopsep=0pt,parsep=0pt,itemsep=2pt,leftmargin=1.4em}
\setlist[enumerate]{topsep=2pt,partopsep=0pt,parsep=0pt,itemsep=2pt,leftmargin=1.6em}
\usepackage{parskip}
\usepackage{fvextra}
\fvset{breaklines=true,breakanywhere=true,fontsize=\small}
\setlength{\parindent}{0pt}

\begin{document}

\begin{center}
  {\LARGE\bfseries Trigonometry}\\[0.35em]
  {\large Igcse Mathematics}
\end{center}
\vspace{0.6em}

This handout covers Topic 6, Trigonometry. Parts marked \textbf{(Extended)} are only tested on the Extended papers; everything else is for both levels. Give angle answers to one decimal place.

\section*{Pythagoras' theorem}

\textbf{Pythagoras' theorem} 勾股定理 links the three sides of a \textbf{right-angled triangle} 直角三角形 (a \textbf{triangle} 三角形 with one $90^{\circ}$ angle). If the longest side (the \textbf{hypotenuse} 斜边, opposite the right angle) is $c$, then

\[
a^{2} + b^{2} = c^{2}.
\]

Use it to find a missing side.

\textbf{Worked example.} A right-angled triangle has a hypotenuse of $13\text{ cm}$ and one short side of $5\text{ cm}$. Find the other short side.

\[
b^{2} = 13^{2} - 5^{2} = 169 - 25 = 144, \qquad b = \sqrt{144} = 12\text{ cm}.
\]

\section*{Trigonometry in right-angled triangles}

Label the sides from the angle you are using: the \textbf{opposite} 对边 (across from the angle), the \textbf{adjacent} 邻边 (next to the angle), and the hypotenuse. The three ratios are the \textbf{sine} 正弦, \textbf{cosine} 余弦 and \textbf{tangent ratio} 正切 (sin, cos, tan):

\[
\sin\theta = \frac{\text{opp}}{\text{hyp}}, \qquad \cos\theta = \frac{\text{adj}}{\text{hyp}}, \qquad \tan\theta = \frac{\text{opp}}{\text{adj}}.
\]

Remember them as \textbf{SOH-CAH-TOA}. To find an angle, use the inverse ($\sin^{-1}$, $\cos^{-1}$, $\tan^{-1}$).

\textbf{Worked example (find a side).} In a right-angled triangle the hypotenuse is $10\text{ cm}$ and the angle is $30^{\circ}$. Find the opposite side.

\[
\text{opp} = 10 \times \sin 30^{\circ} = 10 \times 0.5 = 5\text{ cm}.
\]

\textbf{Worked example (find an angle).} The opposite side is $4\text{ cm}$ and the adjacent side is $3\text{ cm}$.

\[
\tan\theta = \frac{4}{3}, \qquad \theta = \tan^{-1}\!\left(\frac{4}{3}\right) = 53.1^{\circ}.
\]

\section*{Angles of elevation and depression (Extended)}

The \textbf{angle of elevation} 仰角 is the angle up from the horizontal to an object above you. The \textbf{angle of depression} 俯角 is the angle down from the horizontal to an object below you. The shortest distance from a point to a line is the \textbf{perpendicular} 垂直 distance.

\textbf{Worked example.} From a point $50\text{ m}$ from the foot of a tower, the angle of elevation to the top is $40^{\circ}$. Find the height of the tower.

\[
\text{height} = 50 \times \tan 40^{\circ} = 50 \times 0.839 = 42.0\text{ m}.
\]

\section*{Exact trigonometric values (Extended)}

You must know these exact values without a calculator.

\begin{center}
\adjustbox{max width=\linewidth}{%
\begin{tabular}{llllll}
\toprule
\textbf{$x$} & \textbf{$0^{\circ}$} & \textbf{$30^{\circ}$} & \textbf{$45^{\circ}$} & \textbf{$60^{\circ}$} & \textbf{$90^{\circ}$} \\
\midrule
$\sin x$ & $0$ & $\tfrac{1}{2}$ & $\tfrac{\sqrt{2}}{2}$ & $\tfrac{\sqrt{3}}{2}$ & $1$ \\
$\cos x$ & $1$ & $\tfrac{\sqrt{3}}{2}$ & $\tfrac{\sqrt{2}}{2}$ & $\tfrac{1}{2}$ & $0$ \\
$\tan x$ & $0$ & $\tfrac{1}{\sqrt{3}}$ & $1$ & $\sqrt{3}$ & — \\
\bottomrule
\end{tabular}}
\end{center}

\section*{Graphs and trigonometric equations (Extended)}

For $0^{\circ} \leqslant x \leqslant 360^{\circ}$:

\begin{itemize}
\item $y = \sin x$ is a wave starting at $0$, peaking at $90^{\circ}$, back to $0$ at $180^{\circ}$, down to $-1$ at $270^{\circ}$.

\item $y = \cos x$ is the same wave but starting at $1$.

\item $y = \tan x$ rises steeply and repeats every $180^{\circ}$.

\end{itemize}

A \textbf{trigonometric equation} 三角方程 often has more than one answer in this range. Use the graph (or the symmetry of the wave) to find them all.

\textbf{Worked example.} Solve $\sin x = \tfrac{\sqrt{3}}{2}$ for $0^{\circ} \leqslant x \leqslant 360^{\circ}$.

One answer is $x = 60^{\circ}$. The sine wave is also $\tfrac{\sqrt{3}}{2}$ at $180^{\circ} - 60^{\circ} = 120^{\circ}$. So $x = 60^{\circ}$ or $120^{\circ}$.

\textbf{Worked example.} Solve $2\cos x + 1 = 0$ for $0^{\circ} \leqslant x \leqslant 360^{\circ}$.

\[
\cos x = -\tfrac{1}{2} \;\Rightarrow\; x = 120^{\circ} \text{ or } 240^{\circ}.
\]

\section*{The sine and cosine rules (Extended)}

For \textbf{any} triangle (not just right-angled), with sides $a, b, c$ opposite angles $A, B, C$:

\[
\text{sine rule:}\quad \frac{a}{\sin A} = \frac{b}{\sin B} = \frac{c}{\sin C},
\]

\[
\text{cosine rule:}\quad a^{2} = b^{2} + c^{2} - 2bc\cos A.
\]

Use the \textbf{sine rule} 正弦定理 when you have a side and its opposite angle. Use the \textbf{cosine rule} 余弦定理 when you have two sides and the angle between them, or all three sides. With the sine rule, watch for the \textbf{ambiguous case} 两解情况, where an angle could be acute or obtuse.

The area of any triangle is

\[
\text{area} = \tfrac{1}{2}ab\sin C.
\]

\textbf{Worked example.} A triangle has $b = 7\text{ cm}$, $c = 8\text{ cm}$ and the angle between them $A = 40^{\circ}$. Find side $a$.

\[
a^{2} = 7^{2} + 8^{2} - 2(7)(8)\cos 40^{\circ} = 113 - 112 \times 0.766 = 27.2,
\]

\[
a = \sqrt{27.2} = 5.2\text{ cm}.
\]

\section*{Pythagoras and trigonometry in 3D (Extended)}

In three dimensions, find a right-angled triangle inside the solid, then use Pythagoras or trigonometry on it. A common task is the angle between a line and a flat surface (a \textbf{plane} 平面).

\textbf{Worked example.} A box has a base $6\text{ cm}$ by $8\text{ cm}$ and height $5\text{ cm}$. Find the angle between a space diagonal and the base.

First the base diagonal: $\sqrt{6^{2} + 8^{2}} = \sqrt{100} = 10\text{ cm}$. This diagonal and the height form a right-angled triangle, so the angle $\theta$ with the base is

\[
\tan\theta = \frac{5}{10} = 0.5, \qquad \theta = \tan^{-1}(0.5) = 26.6^{\circ}.
\]


\end{document}
